\documentclass[11pt]{article}
\usepackage{amsmath,amssymb}
\title{A legitimate critical background norm for completed quantum packets}\author{}\date{}
\begin{document}\maketitle
\paragraph{Physical object and scope.}
Let $F_R$ be the COMPLETE first quantum coefficient $C_R$ as a function
of the original boundary record, or one complete gauge-covariant
proper-time packet from the longest time completed contact packet calculation. For example a dyadic longest-line packet is
$C(X_{2s},Y_{2s})-C(X_s,Y_s)$, with all vertices, means, ghosts,
Haar and actual background dependence retained. At each fixed box
these objects are boundary-gauge invariant and $DF_R(0)=0$ by global
color invariance. Neither property is asserted for a diagram alone.
We consider their physical second boundary derivative, not the fourth
source derivative of the classical minimum action.

\paragraph{1. Gauge quotient of the actual harmonic tangent.}
Let $h$ be a complete Maxwell-harmonic field extending the actual
boundary tangent. On the FULL original vertex/link graph, including
boundary links, minimize $\|h-D\chi\|_2^2$ over vertex functions
modulo constants. Put $b=h-D\chi$. Then
\[
 D^*b=0\ \hbox{at every graph vertex},\qquad db=dh,
 \qquad (d^*db)_I=0.                                        \tag{1}
\]
This is an honest simultaneous gauge change of the complete boundary
record and the integrated interior, not an independent change of block
frames. Maxwell harmonicity is preserved because $dD=0$.
Differentiating exact boundary gauge covariance, using $DF_R(0)=0$,
gives the two-slot identity
\[
 \boxed{D^2F_R(0)[h,h]=D^2F_R(0)[b,b].}                       \tag{2}
\]
The source curve induced by $\exp(\epsilon\chi)$ may have acceleration;
its contribution here is $DF_R(0)$ applied to that acceleration, hence
zero. This argument cannot be used to discard the corresponding terms
in a fourth-order classical source coefficient. For a single Cartan
color all exponentials commute and the straight source paths already
map exactly; the general quadratic statement also follows by flat
color polarization.

The local-chart cutoff and all gauge/frame factors are transformed
with the ORIGINAL source. This is a fixed-volume coefficient statement;
it supplies no uniform nonlinear chart-tail estimate.

\paragraph{2. Uniform critical norm on the original lattice.}
Reflect $b$ exactly as in the general box magnetic current bound calculation, odd in its own edge direction
with zero normal fold edges and even in the transverse coordinates.
The resulting periodic field $B$ on periods $2(R+1)$ satisfies
\[
 D^*B=0,\qquad \sum_{i,j}\|\delta_jB_i\|_2^2
                       =16\|dh\|_2^2.                      \tag{3}
\]
Every component has zero periodic mean by its odd normal reflection.
The four-dimensional discrete periodic Sobolev inequality consequently
implies
\[
 \boxed{\|b\|_{\ell^4}^2\le C\|dh\|_2^2,
 \quad \|b\|_{\ell^2}\le CR\|dh\|_2,
 \quad \|\nabla B\|_2=4\|dh\|_2.}                          \tag{4}
\]
All constants are independent of $R$ and the boundary data.

For clarity the Sobolev step introduces no lattice-scale constant.
Multilinearly interpolate each periodic component on the unit cells.
Fixed finite-dimensional cell norm equivalence compares its continuum
$L^4$ norm with the lattice $\ell^4$ norm, and bounds its continuum
Dirichlet energy by the sum of squared nearest-neighbor differences.
The interpolant has zero mean. On a torus of side $L$, the elementary
four-dimensional Sobolev inequality is
$\|u\|_4\le C(\|\nabla u\|_2+L^{-1}\|u\|_2)$;
it follows by a cutoff over a fixed number of periodic cells and the
Euclidean four-dimensional Sobolev inequality. Periodic Poincare for
zero mean gives $L^{-1}\|u\|_2\le C\|\nabla u\|_2$.
Both terms have the same scaling, so $C$ is independent of $L$.
Finally each original internal edge is repeated 16 times in the
reflection, while fold edges vanish. This proves(4) with the original
complete curl energy and all boundary multiplicities.

\paragraph{3. Consequence and exact remaining packet estimate.}
For the offset family of the harmonic offset boundary exchange calculation, the constant temporal component is a pure
gauge tangent. Equation(2) removes it in EVERY complete dyadic packet
at this second source order. Thus its classical quartic growth does
not imply a constant-offset growth for $C_R''$.
More generally(4) supplies a valid critical background norm for the
physical quantum Hessian. It is NOT a curl bound on arbitrary tree
fluctuations, nor a license to drop ghost, minimum or boundary partners.

The remaining quantitative problem is the estimate of the FULL signed
packet Hessian on these representatives. Writing $P_{R,s}$ for the
above dyadic packet, one still needs a bound such as
\[
 |D^2P_{R,s}(0)[b,b]|\le c_{R,s}\|dh\|_2^2
\]
with a controlled sum of $c_{R,s}$ over proper-time scales.
No such bound is inferred from(4) alone: the internal loop currents,
short-line contacts, induced exchange and joint coarse covariance
remain. The new valid inequality is the uniform background norm(4)
under the exact physical identity(2). It isolates the correct norm
without assuming an unjustified raw-tree Sobolev estimate.
The displayed derivatives use the physical background parameter. In
$B^{-1}C_R(\epsilon h)$ with canonical $\epsilon=\zeta/\sqrt B$,
the target curvature contributes $B^{-2}D^2C_R(0)[h,h]$; this original
coupling factor is not changed by the gauge quotient.
The first nonlinear coefficient bound, its uniform interacting
remainder and S94 remain open.
\end{document}
